\scanpage{6}
\mspage{5}

commute à $\forall$ changement de base $Y'\longrightarrow Y$.

\margin{\uline{NB} Vrai sans hyp. de propreté si $Y'\longrightarrow Y$ est \uline{lisse}, $f$ qu. cpt. qu. sép., et $F$ premier aux car. rés. de $Y$.}

\deleted{8°) \uline{Ex} Morphismes finis $R^q f_* F=0$ si $q\ne0$}

\[
\left\{\begin{array}{ll}
\uline{\text{Ex}} & R^q f_*(F)_{\bar y}\simeq H^q(X_{\bar y},F_{\bar y})\\
\uline{\text{Ex.}} & X\ \underset{\text{sur}}{\text{prop.}}\ k\text{ sép. clos},\quad K'\text{ ext. sép. close de }k,
\end{array}\right.
\]
d'où $X'=X_{K'}$, $F'=F_{X'}$ et $H^i(X,F)\simeq H^i(X',F')$.

\uline{NB} Vrai sans hyp de propreté si $F$ de torsion premier \deleted{\unsure{à la car.}} car $k$.

8°) \uline{Comparaison avec la cohomologie classique}

$X$ préschéma de \unsure{t.f.} sur $\mathbb C$, alors pour $\forall$ $F$ faisceau de torsion sur $X$,
\[
H^q(X,F)\longrightarrow H^q(X^{\mathrm{top}},F^{\mathrm{top}})
\]
est un \uline{isomorphisme}. Résultat analogue pour un morphisme $f:X\longrightarrow Y$, si $f$ est \deleted{\unsure{propre}} \uline{de type fini}.

\margin{Utilise en plus \unsure{résolution sing. pour l'hyp.} de $f$ \deleted{\unsure{de}} de projectivité.}

\margin{\uline{Gaga} plus complet : O.K. pour \deleted{\unsure{Hom}} $\unsure{\operatorname{Ext}^i|\operatorname{Ext}}$ \unsure{de|les} faisceaux de torsion ; équivalence des catégories des faisceaux constructibles algébriques et faisceaux constructibles analytiques qui sont \uline{de torsion}.}

\deleted{\uline{Remarque} Soit $X$ un schéma propre sur une \unsure{algèbre} analytique \inserted{$A=\mathcal O_{Y,y}$} sur $C$, où $Y$ est un espace analytique sur $C$. Alors $X$ ``provient'' d'un espace analytique $X'$ \uline{propre} au dessus d'un voisinage ouvert $Y'$ de $Y$. \deleted{\unsure{À chaque}} $\unsure{\forall}$ faisceau \inserted{de torsion} constr. $F'$ sur $X'$ ``induit'' un \deleted{\unsure{faisceau}} constr. $F$ sur $X$.}
